\section{Radial exact WKB: add the origin as monodromy data}
\label{sec:radial-wkb}
% BEGIN CURATED SUBJECT INDEX
\index{scattering!Coulomb}
\index{exact WKB}
\index{exact quantization condition}
\index{monodromy}
\index{Langer correction}
\index{Frobenius solution}
% END CURATED SUBJECT INDEX

\calculationroute
  {The global connection algorithm of Stage I and the local Frobenius theory
  of a regular singular point.}
  {Include the origin index and Langer contribution in open-path and
  closed-cycle exact quantization conditions.}
  {Recover oscillator and Coulomb spectra before introducing a
  monodromy-preserving deformation or optimization.}

\subsection{Radial equation and the origin}
\label{subsec:radial-equation}
% BEGIN CURATED SUBJECT INDEX
\index{monodromy}
\index{Langer correction}
\index{Frobenius solution}
\index{radial equation}
% END CURATED SUBJECT INDEX

For a central potential $V(r)$ in three dimensions, write the wave function as
$u_l(r)Y_{lm}(\Omega)/r$.  The reduced radial function obeys
\begin{equation}
  \left[
  -\frac{\hbar^2}{2m}\frac{\rmd^2}{\rmd r^2}
  +V(r)
  +\frac{\hbar^2l(l+1)}{2mr^2}
  \right]u_l(r)
  =Eu_l(r),
  \qquad r>0 .
  \label{eq:radial-schrodinger}
\end{equation}
Here $l\in\naturalnumbers\cup\{0\}$ is the orbital angular momentum.  The
origin is a regular singular point.  Its Frobenius solutions behave as
\begin{equation}
  \vsup{u_l}{reg}(r)\sim r^{l+1},
  \qquad
  \vsup{u_l}{sing}(r)\sim r^{-l} .
  \label{eq:radial-frobenius}
\end{equation}
Physical regularity selects the first solution.

The centrifugal term is proportional to $\hbar^2$, but it cannot be treated
as an innocuous high-order perturbation near $r=0$.  It determines the local
indices and hence the monodromy.  The Langer organization replaces
\begin{equation}
  l(l+1)\longrightarrow\left(l+\frac{1}{2}\right)^2
  \label{eq:langer-replacement}
\end{equation}
inside the leading WKB momentum
\cite{Langer:PhysRev.51.669}.  The additional $1/4$ is the contribution of the
WKB half-density near the regular singularity.

\subsection{Open-path and closed-cycle quantization}
\label{subsec:open-closed}
% BEGIN CURATED SUBJECT INDEX
\index{spectrum!point}
\index{turning point}
\index{exact quantization condition}
\index{monodromy}
\index{Frobenius solution}
% END CURATED SUBJECT INDEX

There are two natural formulations.  In the open-path formulation, transport
a decaying solution from infinity to the origin and require the coefficient
of $\vsup{u_l}{sing}$ to vanish.  In the closed-cycle formulation, continue
the same solution around a contour that encloses the origin and relevant
turning points, then demand a compatible monodromy.  The two conditions are
equivalent when all connection matrices and the local origin monodromy are
included
\cite{AokiKawaiTakei1991RIMS853,Kawai:2005,
AokiIwakiTakahashi2019LoopType,Morikawa:2025ezl}.

The equivalence can be stated schematically as
\begin{equation}
  \vsub{C}{sing}(E,\hbar)=0
  \quad\Longleftrightarrow\quad
  \det\left[M_{\Gamma}(E,\hbar)-\identity\right]=0 ,
  \label{eq:open-closed-equivalence}
\end{equation}
where $\vsub{C}{sing}$ is the open-path coefficient of the singular
Frobenius solution and $M_{\Gamma}$ is the ordered monodromy matrix around a
closed contour $\Gamma$.  The determinant notation is schematic because a
specific quantization condition may select one eigenvalue or one matrix
element rather than the full determinant.  What is invariant is the vanishing
of the undesired component.

The change of variables
\begin{equation}
  r=\rme^x,
  \qquad x\in(-\infty,\infty),
  \label{eq:radial-exponential-map}
\end{equation}
makes the equivalence intuitive.  The punctured origin becomes the boundary
$x\to-\infty$.  A small-circle monodromy in the $r$ plane becomes a boundary
condition at the left end of the $x$ line, while $r\to\infty$ becomes
$x\to+\infty$.  The radial spectral problem is thereby converted into a
two-ended connection problem.

\subsection{Harmonic-oscillator benchmark}
\label{subsec:radial-oscillator}
% BEGIN CURATED SUBJECT INDEX
\index{topology}
\index{turning point}
\index{monodromy}
\index{Langer correction}
% END CURATED SUBJECT INDEX

For
\begin{equation}
  V(r)=\frac{1}{2}m\Omega^2r^2,
  \qquad \Omega>0,
  \label{eq:radial-oscillator-potential}
\end{equation}
the exact spectrum is
\begin{equation}
  E_{n_rl}
  =\hbar\Omega
  \left(2n_r+l+\frac{3}{2}\right),
  \qquad n_r=0,1,2,\ldots .
  \label{eq:radial-oscillator-spectrum}
\end{equation}
A connection using only the two ordinary turning points misses the
$l$ dependence.  Adding the origin monodromy, or equivalently using the Langer
momentum from the start, restores the exact result.  This is a clean example
of an order-$\hbar^2$ term carrying order-one topological data.

\subsection{Coulomb benchmark}
\label{subsec:radial-coulomb}
% BEGIN CURATED SUBJECT INDEX
\index{scattering!Coulomb}
\index{turning point}
\index{Voros symbol}
% END CURATED SUBJECT INDEX

For an attractive Coulomb potential
\begin{equation}
  V(r)=-\frac{g}{r},
  \qquad g>0,
  \label{eq:coulomb-potential}
\end{equation}
the bound-state spectrum is
\begin{equation}
  E_{n_rl}
  =-\frac{mg^2}{2\hbar^2(n_r+l+1)^2} .
  \label{eq:coulomb-spectrum}
\end{equation}
Again the origin contribution supplies the $l+1/2$ phase, while the ordinary
turning point supplies the remaining Maslov half-unit.  The total principal
quantum number is $n=n_r+l+1$.  Exact evaluation of the closed Voros period
shows that the open and closed quantization paths agree once the regular
singularity is treated consistently
\cite{Morikawa:2025ezl}.

\subsection{Monodromy renormalization}
\label{subsec:monodromy-rg}
% BEGIN CURATED SUBJECT INDEX
\index{scattering!Coulomb}
\index{monodromy}
\index{renormalization group}
% END CURATED SUBJECT INDEX

Local WKB integrals around a singular point can contain regulator-dependent
pieces.  Introduce a small radius $\mu>0$ and a circle $C_0(\mu)$ around the
origin.  A useful coupling that keeps the origin index at leading order is
\begin{equation}
  \lambda=\hbar\left(l+\frac{1}{2}\right) .
  \label{eq:monodromy-coupling}
\end{equation}
After the WKB half-density is included, fix the physical small-circle phase
to
\begin{equation}
  \vsup{\mathcal M_0}{phys}
  =\rme^{\rmi\pi(2l+1)} .
  \label{eq:physical-origin-monodromy}
\end{equation}
Define a counter-Voros factor $\vsub{Z}{mon}(\mu)$ by
\begin{equation}
  \vsup{\mathcal M_0}{phys}
  =\frac{\vsup{\mathcal M_0}{bare}(\mu)}
  {\vsub{Z}{mon}(\mu)} .
  \label{eq:monodromy-Z}
\end{equation}
Cutoff independence is summarized by an RG-type equation
\begin{equation}
  \mu\frac{\rmd}{\rmd\mu}
  \log \vsub{Z}{mon}(\mu)
  =\vsub{\gamma}{mon}(E,\hbar,\lambda) .
  \label{eq:monodromy-rge}
\end{equation}
The function $\vsub{\gamma}{mon}$ is a monodromy anomalous dimension.  A
closed period may then be split into large and small circles with local
counterterms chosen so that the sum is finite and independent of $\mu$.

Equations~\eqref{eq:monodromy-Z} and \eqref{eq:monodromy-rge} should be read
with care.  They organize regulator independence of connection data; they are
not automatically a Wilsonian coarse graining of spatial degrees of freedom.
For the harmonic oscillator and Coulomb benchmarks, the proposed scheme has
$\vsub{\gamma}{mon}=0$, corresponding to fixed monodromy data and the
termination of higher closed-period corrections in the chosen organization
\cite{Morikawa:2025ezl}.  For a generic potential, establishing the existence,
scheme dependence, and physical content of a nonzero
$\vsub{\gamma}{mon}$ remains an open mathematical problem.

\subsection{Optimized perturbation around monodromy-preserving seeds}
\label{subsec:optimized-wkb}
% BEGIN CURATED SUBJECT INDEX
\index{scattering!Coulomb}
\index{Borel summation}
\index{monodromy}
\index{optimized perturbation theory}
\index{variational perturbation theory}
% END CURATED SUBJECT INDEX

For a non-solvable radial potential, choose a solvable seed $Q^{(0)}$ that has
the same local monodromy and introduce an interpolation parameter $\delta$:
\begin{equation}
  Q_{\delta}(r,E)
  =Q^{(0)}(r,E;\bm a)
  +\delta\left[Q(r,E)-Q^{(0)}(r,E;\bm a)\right] .
  \label{eq:delta-interpolation}
\end{equation}
The vector $\bm a$ contains variational scales.  Expand the Borel-resummed
Voros data in $\delta$, set $\delta=1$, and choose $\bm a$ by minimal
sensitivity or fastest apparent convergence.  Because the seed already
contains the correct origin index, the expansion does not repeatedly rebuild
the most singular local information.  This is an exact-WKB version of
optimized or variational perturbation theory
\cite{Stevenson:1981vj,Duncan:1992ba,Buckley:1992pc,
Kleinert:1995hc,Zinn-Justin:2010zzb}.

The method is promising for Cornell, screened Coulomb, and other radial
potentials, but three checks are indispensable: independence of the auxiliary
parameter at increasing order, compatibility with lateral Borel summation,
and preservation of the physical origin boundary condition.  Variational
  stationarity alone does not prove that a selected branch is the physical
  resurgent continuation.
